\newif\ifglobal

%%%%%%%%%%%%%%%%%%%%%%%
%%% Compile to view %%%
%%%%%%%%%%%%%%%%%%%%%%%

%%%%%%%%%%%%%%%%%%%%%%%%%%%%%%%%%%%%%%%%%%%%%%%%%%%
%%% Readme:                                     %%%
%%% https://www.overleaf.com/read/bwdjvfjksvjy/ %%%
%%%%%%%%%%%%%%%%%%%%%%%%%%%%%%%%%%%%%%%%%%%%%%%%%%%

% >>> M?: marks a module setting number or important notice

% >>> M4: Set {./relative-path} to external files for this module <<<
\def \modulefiles {./16-ECC}
% To add an external file, use:
% (Replace '---' with actual filename)
% '\input{\modulefiles/tables/---.tex}' for tables
% '\includegraphics{\modulefiles/figures/---}' for figures
% '\subfile{\modulefiles/---}' for register files


% >>> M1: This module is in global project? <<<
% 'YES' - uncomment; 'NO' - comment out
\globaltrue

\ifglobal
    \documentclass[main\_\_EN.tex]{subfiles}

\else
    \documentclass[openany, 10pt]{book}

    % >>> M2: In ./00-shared/preamble-trm-module, also find
    %         settings for visibility of todo notes, watermark <<<
    \usepackage{preamble-trm-module}

    % >>> M3: Temporary labels in English,
    %         you can add your own labels there <<<
    \myexternaldocument{./00-shared/temp-labels-en}

\fi %\ifglobal


\setcounter{secnumdepth}{4}


% >>> M5: If this module needs any updates in
% './00-shared/preamble-trm-module.sty'
% (include additional package, etc.), contact your module PM
% or create the file './00-shared/changelog.tex' make the
% necessary updates yourself and document them in this file <<<


\begin{document}


% Sets language into which pre-defined bits of text are translated
\selectlanguage{english}

% Do not delete this block
\ifglobal
% Add the following front matter if
% this module is outside of global project
\else
    \InputIfFileExists{readme}{}{}
    {\let\clearpage\relax \listoftodos}
    {\let\clearpage\relax \listofdonetodos}
    \tableofcontents
    %\listoftables
    %\listoffigures
\fi


%%%%%%%%%%%%%%%%%%%%%%%%%
%%% TRM Module Begins %%%
%%%%%%%%%%%%%%%%%%%%%%%%%


\hypertarget{ecc}{}
\chapter{ECC Hardware Accelerator (ECC)}
\label{mod:ecc}


\section{Introduction}

Elliptic Curve Cryptography (ECC) is an approach to public-key cryptography based on the algebraic structure of elliptic curves. ECC allows smaller keys compared to RSA cryptography while providing equivalent security.

\chipname{}'s ECC Accelerator can complete various calculation based on different elliptic curves, thus accelerating ECC algorithm and ECC-derived algorithms (such as ECDSA).

\section{Features}
\chipname{}'s ECC Accelerator supports:
\begin{itemize}
    \item Two different elliptic curves, namely P-192 and P-256 defined in \href{https://nvlpubs.nist.gov/nistpubs/FIPS/NIST.FIPS.186-5.pdf}{FIPS 186-5}
    \item Seven working modes
    \item Interrupt upon completion of calculation
\end{itemize}

\section{Terminology}
To better illustrate the ECC accelerator, we will first introduce the terminology used in this chapter.

\subsection{ECC Basics}

\subsubsection{Elliptic Curve and Points on the Curves}
The ECC algorithm is based on elliptic curves over prime fields, which can be represented as:
\[
{y}^{2} = {x}^{3} + {ax} + b \ \textrm{mod}\ p
\]
where,
\begin{itemize}
    \item $p$ is a prime number.
    \item $a$ and $b$ are two non-negative integers smaller than $p$.
    \item $(x, y)$ is a point on the curve satisfying the representation.
\end{itemize}

\subsubsection{Affine Coordinates and Jacobian Coordinates}

An elliptic curve can be represented as below:

\begin{itemize}
    \item In affine coordinates:
    \[
    {y}^{2} = {x}^{3} + {ax} + b \ \textrm{mod}\ p
    \]
    \item In a Jacobian coordinates:
    \[
    {Y}^{2} = {X}^{3} + {aXZ}^{4} + {bZ}^{6} \ \textrm{mod}\ p
    \]
\end{itemize}

To convert affine coordinates $(x,y)$ to/from Jacobian coordinates $(X,Y,Z)$:

\begin{itemize}
    \item From Affine to Jacobian coordinates
    \[
    x = X / {Z}^2 \ \textrm{mod}\ p
    \]
    \[
    y = Y / {Z}^3 \ \textrm{mod}\ p
    \]
    \item From Jacobian to affine coordinates
    \[ X = x \]
    \[ Y = y \]
    \[ Z = 1\]
\end{itemize}

\subsection{ECC Definitions}

\subsubsection{Memory Blocks}

ECC's memory blocks store input date and output data of the ECC operation.

\begin{table}[H]
    \centering
    \caption{ECC Accelerator Memory Blocks} \label{tab:ecc-memory}
    \begin{threeparttable}
    \begin{tabular}{ | l | l | l | l | l | l | }
    \hline
        \rowcolor{lightgray}
        \textbf{Memory} & \textbf{Size (byte)}  &  \textbf{Starting Address}\tnote{*} &  \textbf{Ending Address} \tnote{*} & \textbf{Access}   \\ \hline
        ECC\_MULT\_Mem\_k    & 32  & 0x100 & 0x11F & R/W \\ \hline
        ECC\_MULT\_Mem\_Px   & 32  & 0x120 & 0x13F & R/W \\ \hline
        ECC\_MULT\_Mem\_Py   & 32  & 0x140 & 0x15F & R/W \\ \hline
    \end{tabular}
    \begin{tablenotes}
        \item[*] Address offset relative to ECC accelerator base address provided in Table \ref{tab:sysmem-base-address} in Chapter \ref{mod:sysmem} \textit{\nameref{mod:sysmem}}.
    \end{tablenotes}
    \end{threeparttable}
\end{table}

\subsubsection{Data and Data Block}

\chipname{}'s ECC operates on data of 256 bits. This data ($D[255:0]$) can be divided into eight 32-bit data blocks $D[n][31:0] (n=0, 1, \cdots, 7)$. To be specific:
\[
D[255:0] = {D[7][31:0],D[6][31:0],D[5][31:0],
D[4][31:0],D[3][31:0],D[2][31:0],D[1][31:0],D[0][31:0]}
\]

\subsubsection{Write Data}

Write data means writing data to an ECC memory block and using this data as the input to the ECC algorithm. To be more specific, write data to an ECC memory block means write $D[n][31:0] (n=0, 1, \cdots, 7)$ to the ``starting address of this ECC memory block +$
4 \times n$" successively:

\begin{itemize}
    \item write $D[0]$ to ``starting address"
    \item write $D[1]$ to ``starting address + 4"
    \item $\cdots$
    \item write $D[7]$ to ``starting address + 28"
\end{itemize}

\begin{tiplisting}
When the key size of 192 bits is used, you need to append 0 before 192 bits of data and write 256 bits of data.
\end{tiplisting}

\subsubsection{Read Data}

Read data means reading data from the starting address and using this data as the output from the ECC algorithm. To be more specific, read data from an ECC memory block means read $D[n][31:0] (n=0, 1, \cdots, 7)$ from the ``starting address of this ECC memory block +
$4 \times n$" successively:
\begin{itemize}
    \item read $D[0]$ from ``starting address"
    \item read $D[1]$ from ``starting address + 4"
    \item $\cdots$
    \item read $D[7]$ from ``starting address + 28"
\end{itemize}

\begin{tiplisting}
When the key size of 192 bits is used, only read 192 bit (6 blocks) of data.
\end{tiplisting}


\subsubsection{Standard Calculation and Jacobian Calculation}

\chipname{}'s ECC performs Base Point Calculation (including Base Point Verification and Base Point Multiplication) using the affine coordinates and Jacobian Calculation (including Jacobian Point Verification and Jacobian Point Multiplication) using the Jacobian coordinates.

\section{Function Description}
\label{sec:ecc-mode-config}
\subsection{Key Size}

\chipname{}'s ECC supports acceleration based on two key sizes (corresponding to two different elliptic curves). By configuring \hyperref[fielddesc:ECCMULTKEYLENGTH]{ECC\_MULT\_KEY\_LENGTH} field, users can choose desired key size. Details can be seen in Table \ref{tab:ecc-key-length} below.

\begin{table}[H]
    \centering
    \caption{Choose ECC Accelerator Key Size}
    \label{tab:ecc-key-length}
    \begin{threeparttable}
    \begin{tabular}{|p{4.2cm}|p{3.6cm}|}
    \hline
        \rowcolor{lightgray}
        \textbf{\hyperref[fielddesc:ECCMULTKEYLENGTH]{ECC\_MULT\_KEY\_LENGTH}} & \textbf{Elliptic Curves} \\ \hline
        1'b0 & FIPS P-192  \\\hline
        1'b1 & FIPS P-256  \\\hline
    \end{tabular}
    \begin{tablenotes}
        \item[1] See definition of FIPS P-192 and P-256 in \href{https://nvlpubs.nist.gov/nistpubs/FIPS/NIST.FIPS.186-5.pdf}{FIPS 186-5}.
    \end{tablenotes}
    \end{threeparttable}
\end{table}

\subsection{Working Modes}

\chipname{}'s ECC accelerator supports 7 working modes based on two elliptic curves described in the above section. By configuring \hyperref[fielddesc:ECCMULTWORKMODE]{ECC\_MULT\_WORK\_MODE} field, users can choose desired working mode. Details can be seen in Table \ref{tab:ecc-work-mode}.


\begin{table}[H]
    \centering
    \caption{ECC Accelerator's Working Modes}
    \label{tab:ecc-work-mode}
    \begin{tabular}{|l|L{4cm}|l|L{4cm}|}
    \hline
        \rowcolor{lightgray}
        \textbf{\hyperref[fielddesc:ECCMULTWORKMODE]{ECC\_MULT\_WORK\_MODE}} & \textbf{Working Modes} & \textbf{\hyperref[fielddesc:ECCMULTWORKMODE]{ECC\_MULT\_WORK\_MODE}} & \textbf{Working Modes} \\ \hline
        3'd0 & Point Multi Mode       & 3'd4 & Jacobian Point Multi             \\\hline
        3'd1 & Division Mode         & 3'd5 & \textbf{Reserved}  \\\hline
        3'd2 & Point Verif   & 3'd6 & Jacobian Point Verif          \\\hline
        3'd3 & Point Verif + Multi  & 3'd7 & Point Verification + Jacobian Multi \\\hline
    \end{tabular}
\end{table}

Detailed description about each working modes is provided in the following sections.

\subsubsection{Base Point Multiplication (Point Multi Mode)}

Base Point Multiplication can be represented as:
\[
(Q_x,Q_y) = k \cdot (P_x,P_y)
\]
where,
\begin{itemize}
    \item Input: $P_x$, $P_y$, and $k$  are stored in \hyperref[tab:ecc-memory]{ECC\_MULT\_Mem\_Px}, \hyperref[tab:ecc-memory]{ECC\_MULT\_Mem\_Py}, and \hyperref[tab:ecc-memory]{ECC\_MULT\_Mem\_k} respectively.
    \item Output: $Q_x$ and $Q_y$ are stored in \hyperref[tab:ecc-memory]{ECC\_MULT\_Mem\_Px} and \hyperref[tab:ecc-memory]{ECC\_MULT\_Mem\_Py} respectively.
\end{itemize}

\subsubsection{Finite Field Division (Division Mode)}

Finite Field Division can be represented as:
\[
\textrm{Result} = P_y \cdot k^{-1}
\]
where,
\begin{itemize}
    \item Input: $P_y$ and $k$ are stored in \hyperref[tab:ecc-memory]{ECC\_MULT\_Mem\_Py} and \hyperref[tab:ecc-memory]{ECC\_MULT\_Mem\_k}.
    \item Output: Result is stored in \hyperref[tab:ecc-memory]{ECC\_MULT\_Mem\_Py}.
\end{itemize}

\subsubsection{Base Point Verification (Point Verif Mode)}

Base Point Verification can be used to verify if a point  $(P_x,P_y)$ is on a selected elliptic curve.

\begin{itemize}
    \item Input: $P_x$ and $P_y$ are stored in \hyperref[tab:ecc-memory]{ECC\_MULT\_Mem\_Px} and \hyperref[tab:ecc-memory]{ECC\_MULT\_Mem\_Py}, respectively.
    \item Output: verification result is stored in \hyperref[fielddesc:ECCMULTVERIFICATIONRESULT]{ECC\_MULT\_VERIFICATION\_RESULT} field.
\end{itemize}

\subsubsection{Base Point Verification + Base Point Multiplication (Point Verif + Multi Mode)}

In this working mode, ECC first verifies if Point $(P_x,P_y)$ is on the selected elliptic curve or not. If yes, then perform the multiplication:
\[
(Q_x,Q_y) = k \cdot (P_x,P_y)
\]
where,
\begin{itemize}
    \item Input: $P_x$, $P_y$ and $k$ are stored at \hyperref[tab:ecc-memory]{ECC\_MULT\_Mem\_Px}, \hyperref[tab:ecc-memory]{ECC\_MULT\_Mem\_Py}, and \hyperref[tab:ecc-memory]{ECC\_MULT\_Mem\_k} respectively.
    \item Output:
    \begin{itemize}
        \item verification result is stored in \hyperref[fielddesc:ECCMULTVERIFICATIONRESULT]{ECC\_MULT\_VERIFICATION\_RESULT} field.
        \item $Q_x$ and $Q_y$ are stored in \hyperref[tab:ecc-memory]{ECC\_MULT\_Mem\_Px} and \hyperref[tab:ecc-memory]{ECC\_MULT\_Mem\_Py} respectively.
    \end{itemize}
\end{itemize}

\subsubsection{Jacobian Point Multiplication (Jacobian Point Multi Mode)}

Jacobian Point Multiplication can be represented as:
\[
(Q_x,Q_y,Q_z) = k \cdot (P_x,P_y,1)
\]
where,
\begin{itemize}
    \item $(Q_x,Q_y,Q_z)$ is a Jacobian point on the selected elliptic curve.
    \item 1 in the point's Jacobian coordinates is auto completed by hardware.
    \item Input: $P_x$, $P_y$ and $k$ are stored in \hyperref[tab:ecc-memory]{ECC\_MULT\_Mem\_Px}, \hyperref[tab:ecc-memory]{ECC\_MULT\_Mem\_Py}, and \hyperref[tab:ecc-memory]{ECC\_MULT\_Mem\_k} respectively.
    \item Output: $Q_x$, $Q_y$, and $Q_z$ are stored in \hyperref[tab:ecc-memory]{ECC\_MULT\_Mem\_Px}, \hyperref[tab:ecc-memory]{ECC\_MULT\_Mem\_Py}, and \hyperref[tab:ecc-memory]{ECC\_MULT\_Mem\_k}, respectively.
\end{itemize}

\subsubsection{Jacobian Point Verification (Jacobian Point Verif Mode)}

Jacobian Point Verification can be used to verify if a point $(Q_x,Q_y,Q_z)$ is on a selected elliptic curve.
\begin{itemize}
    \item $(Q_x,Q_y,Q_z)$ is the point in Jacobian Coordinates.
    \item Input: $Q_x$, $Q_y$, and $Q_z$ are stored in \hyperref[tab:ecc-memory]{ECC\_MULT\_Mem\_Px}, \hyperref[tab:ecc-memory]{ECC\_MULT\_Mem\_Py}, and \hyperref[tab:ecc-memory]{ECC\_MULT\_Mem\_k}, respectively.
    \item Output: verification result is stored in \hyperref[fielddesc:ECCMULTVERIFICATIONRESULT]{ECC\_MULT\_VERIFICATION\_RESULT} field.
\end{itemize}

\subsubsection{Base Point Verification + Jacobian Point Multiplication (Point Verif + Jacobian Point Multi Mode)}

In this working mode, ECC first verifies if Point $(P_x,P_y)$ is on the selected elliptic curve or not. If yes, then perform the multiplication:
\[
(Q_x,Q_y,Q_z) = k \cdot (P_x,P_y,1)
\]
where,
\begin{itemize}
    \item $(Q_x,Q_y,Q_z)$ is a Jacobian point on the selected elliptic curve.
    \item 1 in the point's Jacobian coordinates is auto completed by hardware.
    \item Input: $P_x$, $P_y$, and $k$ are stored in \hyperref[tab:ecc-memory]{ECC\_MULT\_Mem\_Px}, \hyperref[tab:ecc-memory]{ECC\_MULT\_Mem\_Py}, and \hyperref[tab:ecc-memory]{ECC\_MULT\_Mem\_k}.
    \item Output:
    \begin{itemize}
        \item verification result is stored in  \hyperref[fielddesc:ECCMULTVERIFICATIONRESULT]{ECC\_MULT\_VERIFICATION\_RESULT} field.
        \item $Q_x$, $Q_y$, and $Q_z$ are stored in \hyperref[tab:ecc-memory]{ECC\_MULT\_Mem\_Px}, \hyperref[tab:ecc-memory]{ECC\_MULT\_Mem\_Py}, and \hyperref[tab:ecc-memory]{ECC\_MULT\_Mem\_k}.
    \end{itemize}
\end{itemize}

\section{Clocks and Resets}
\chipname{}'s ECC only has one clock module (crypo\_ecc\_clk) and one reset module (crypto\_ecc\_rst). Users should enable the ECC clock and disable the ECC reset before starting the ECC accelerator. For details on how to configure the ECC clock and reset, please refer to Chapter \ref{mod:resclk} \textit{\nameref{mod:resclk}}.


\section{Interrupts}
\label{sec:ecc-interrupt}
\chipname{}'s ECC accelerator can generate one interrupt signal ECC\_INTR and send it to \hyperref[mod:intmtrx]{Interrupt Matrix}.

\begin{tiplisting}
Each interrupt signal is generated by any of its interrupts: any of its interrupt triggered can generate the interrupt signal.
\end{tiplisting}

ECC\_INTR has only one interrupt \label{int:ECCMULTCALCDONEINT}ECC\_MULT\_CALC\_DONE\_INT, which is triggered on the completion of an ECC computation.


This ECC interrupt ECC\_MULT\_CALC\_DONE\_INT is configured by the following registers:
\begin{itemize}
    \item \hyperref[fielddesc:ECCMULTCALCDONEINTRAW]{ECC\_MULT\_CALC\_DONE\_INT\_RAW}: stores the raw interrupt of \hyperref[int:ECCMULTCALCDONEINT]{ECC\_MULT\_CALC\_DONE\_INT}.
    \item \hyperref[fielddesc:ECCMULTCALCDONEINTST]{ECC\_MULT\_CALC\_DONE\_INT\_ST}: indicates the status of the \hyperref[int:ECCMULTCALCDONEINT]{ECC\_MULT\_CALC\_DONE\_INT} interrupt. This field is generated by enabling/disabling  \hyperref[fielddesc:ECCMULTCALCDONEINTRAW]{ECC\_MULT\_CALC\_DONE\_INT\_RAW} field via \hyperref[fielddesc:ECCMULTCALCDONEINTENA]{ECC\_MULT\_CALC\_DONE\_INT\_ENA}.
    \item \hyperref[fielddesc:ECCMULTCALCDONEINTENA]{ECC\_MULT\_CALC\_DONE\_INT\_ENA}: enables/disables the \hyperref[int:ECCMULTCALCDONEINT]{ECC\_MULT\_CALC\_DONE\_INT} interrupt.
    \item \hyperref[fielddesc:ECCMULTCALCDONEINTCLR]{ECC\_MULT\_CALC\_DONE\_INT\_CLR}: set this bit to clear the \hyperref[int:ECCMULTCALCDONEINT]{ECC\_MULT\_CALC\_DONE\_INT} interrupt status. By setting this bit to 1, fields \hyperref[fielddesc:ECCMULTCALCDONEINTRAW]{ECC\_MULT\_CALC\_DONE\_INT\_RAW} and  \hyperref[fielddesc:ECCMULTCALCDONEINTST]{ECC\_MULT\_CALC\_DONE\_INT\_ST} will be cleared.
\end{itemize}

\section{Programming Procedures}

The programming procedures for configuring ECC are described below:

\begin{enumerate}
    \item Configure the ECC clock and reset.
    \item Choose key size and working mode as described in Section \ref{sec:ecc-mode-config}.
    \item Enable \hyperref[int:ECCMULTCALCDONEINT]{ECC\_MULT\_CALC\_DONE\_INT} interrupt as described in Section \ref{sec:ecc-interrupt}.
    \item Set \hyperref[fielddesc:ECCMULTSTART]{ECC\_MULT\_START} field to start ECC calculation.
    \item Wait for the \hyperref[int:ECCMULTCALCDONEINT]{ECC\_MULT\_CALC\_DONE\_INT} interrupt, which indicates the completion of the ECC calculation.
    \item Check the result as described in Section \ref{sec:ecc-mode-config}.
\end{enumerate}

\clearpage
\hypertarget{ecc-reg-summ}{}
\section{Register Summary}
\label{sec:ecc-reg-summ}

The addresses in this section are relative to ECC accelerator base address provided in Table \ref{tab:sysmem-base-address} in Chapter \ref{mod:sysmem} \textit{\nameref{mod:sysmem}}.

The abbreviations given in Column \textbf{Access} are explained in Section \hyperref[glossary-access-types]{\textit{Access Types for Registers}}.

\subfile{\modulefiles/16-ECC-reg-summ__EN}


%%%%%%%%%%%%%%%%%
%%% Registers %%%
%%%%%%%%%%%%%%%%%

% >>> Update the sample text below and insert
%      the info on registers into the subfile <<<

\clearpage
\section{Registers}

The addresses in this section are relative to ECC accelerator base address provided in Table \ref{tab:sysmem-base-address} in Chapter \ref{mod:sysmem} \textit{\nameref{mod:sysmem}}.

\subfile{\modulefiles/16-ECC-reg__EN}


\end{document}
